\chapter{Espaces vectoriels normés}\label{ch:b2:nvs}

Lorsque l'\omterm{def:b2:metric:def}{espace métrique} est un espace vectoriel et que la distance
provient d'une \omterm{def:b2:nvs:norm}{norme}, la topologie et l'\omterm{def:b2:structures:algebra}{algèbre} linéaire commencent à
interagir~: les applications linéaires sont \omterm{def:b2:metric:continuity}{continues} exactement
lorsqu'elles sont bornées sur la boule unité, la dimension finie force
toutes les \omterm{def:b2:nvs:norm}{normes} à coïncider, et la complétude transforme les séries
absolument convergentes en séries convergentes. La distinction
dimension finie/infinie --- cristallisée dans le théorème de Riesz ---
est la leçon la plus profonde du chapitre.

Dans tout le chapitre, $E, F$ sont des espaces vectoriels sur $K = \R$
ou $\C$.

\section{Normes}

\begin{definition}\label{def:b2:nvs:norm}
Une \emph{norme}\index{norme} sur $E$ est une application $\norm{\,\cdot\,} \colon E
\to \R_+$ vérifiant, pour tous $x, y \in E$, $\lambda \in K$~:
\[
\norm x = 0 \iff x = 0,
\qquad
\norm{\lambda x} = \abs\lambda\,\norm x,
\qquad
\norm{x + y} \leq \norm x + \norm y .
\]
Alors $d(x, y) = \norm{x - y}$ est une distance, et tout le
\cref{ch:b2:metric} s'applique. L'inégalité triangulaire inversée
$\bigl|\norm x - \norm y\bigr| \leq \norm{x - y}$ rend la norme
elle-même $1$-lipschitzienne~; l'addition et la multiplication par un
scalaire sont \omterm{def:b2:metric:continuity}{continues} (estimations $\norm{(x + y) - (x' + y')} \leq
\norm{x - x'} + \norm{y - y'}$, etc.).
\end{definition}

\begin{example}\label{ex:b2:nvs:examples}
Sur $K^n$~:
\[
\norm{x}_1 = \sum_i \abs{x_i},
\qquad
\norm{x}_2 = \Bigl(\sum_i \abs{x_i}^2\Bigr)^{1/2},
\qquad
\norm{x}_\infty = \max_i \abs{x_i}
\]
($\norm\cdot_2$ est une \omterm{def:b2:nvs:norm}{norme} d'après Cauchy--Schwarz, volume de
première année). Sur $C(\intcc{a}{b})$~:
\[
\norm f_\infty = \sup \abs f,
\qquad
\norm f_1 = \int_a^b \abs f,
\qquad
\norm f_2 = \Bigl(\int_a^b \abs f^2\Bigr)^{1/2},
\]
les deux dernières étant des \omterm{def:b2:nvs:norm}{normes} grâce à la stricte positivité de
l'intégrale et à l'inégalité de Cauchy--Schwarz intégrale (volume de
première année). Sur les matrices~: n'importe quelle \omterm{def:b2:nvs:norm}{norme} sur
$\mathcal M_n(K) \simeq K^{n^2}$~; les \omterm{thm:b2:nvs:continuouslinear}{normes d'opérateur} ci-dessous
sont les plus importantes structurellement.
\end{example}

\begin{definition}[Normes équivalentes]\label{def:b2:nvs:equivalent}
Deux \omterm{def:b2:nvs:norm}{normes} $N_1, N_2$ sur $E$ sont \emph{équivalentes}\index{normes
équivalentes} lorsqu'il existe des constantes $c, C > 0$ telles que
\[
c\,N_1 \leq N_2 \leq C\, N_1 .
\]
Deux \omterm{def:b2:nvs:norm}{normes} équivalentes ont les mêmes \omterm{def:b2:metric:topology}{ouverts}, les mêmes suites
convergentes et de Cauchy, les mêmes parties \omterm{def:b2:metric:compact}{compactes} et complètes~:
la même analyse.
\end{definition}

\begin{example}[Non-équivalence en dimension infinie]\label{ex:b2:nvs:nonequivalent}
Sur $C(\intcc{0}{1})$~: $\norm f_1 \leq \norm f_\infty$ toujours, mais
aucune borne inverse ne tient~: $f_n(x) = x^n$ vérifie $\norm{f_n}_\infty = 1$
et $\norm{f_n}_1 = \frac{1}{n+1} \to 0$. Ainsi $f_n \to 0$ pour
$\norm\cdot_1$ mais pas pour $\norm\cdot_\infty$~: les deux \omterm{def:b2:nvs:norm}{normes}
sont en désaccord sur la convergence elle-même.
\end{example}

\begin{example}[Constantes explicites en dimension $n$]\label{ex:b2:nvs:constants}
Sur $K^n$, les trois \omterm{def:b2:nvs:norm}{normes} classiques sont \omterm{def:b2:nvs:equivalent}{équivalentes} avec des
constantes optimales~:
\[
\norm x_\infty \leq \norm x_2 \leq \norm x_1
\leq \sqrt n\,\norm x_2 \leq n\,\norm x_\infty ,
\]
la borne médiane $\norm x_1 \leq \sqrt n\norm x_2$ provenant de
Cauchy--Schwarz contre le vecteur constant égal à $1$. Vecteurs
extrémaux~: $e_1$ rend les deux premières inégalités des égalités,
$(1, 1, \dots, 1)$ les deux dernières. La dimension $n$ figure
visiblement dans les constantes --- la graine quantitative de l'échec
en dimension infinie~: lorsque $n \to \infty$, aucune constante
uniforme ne survit, ce qui est exactement ce que l'\cref{ex:b2:nvs:nonequivalent}
exhibe sur les espaces de fonctions.
\end{example}

\begin{omfigure}
\begin{tikzpicture}[scale=1.5]
  \draw[->] (-1.45,0) -- (1.5,0) node[below, font=\small] {$x$};
  \draw[->] (0,-1.45) -- (0,1.5) node[left, font=\small] {$y$};
  \draw[omProp, very thick]
    (1,1) -- (-1,1) -- (-1,-1) -- (1,-1) -- cycle;
  \draw[omThm, very thick] (0,0) circle (1);
  \draw[omDef, very thick]
    (1,0) -- (0,1) -- (-1,0) -- (0,-1) -- cycle;
  \node[omDef, font=\small] at (0.36,0.28)
    {$\norm\cdot_1$};
  \node[omThm, font=\small] at (0.55,0.92)
    {$\norm\cdot_2$};
  \node[omProp, font=\small] at (1.24,1.13)
    {$\norm\cdot_\infty$};
\end{tikzpicture}

{\small Les boules unités des trois \omterm{def:b2:nvs:norm}{normes} classiques de $\R^2$,
emboîtées comme le dictent les inégalités de l'\cref{ex:b2:nvs:constants}~:
plus la boule est petite, plus la \omterm{def:b2:nvs:norm}{norme} est grande. La rotondité
compte~: les côtés plats du losange et du carré sont exactement les
défauts de stricte convexité exploités dans le problème du week-end du
\cref{ch:b2:realfun} et dans celui de ce chapitre (question 4).}
\end{omfigure}

\section{Applications linéaires continues}

\begin{theorem}[Caractérisation]\label{thm:b2:nvs:continuouslinear}
Pour une application linéaire $u \colon E \to F$ entre espaces normés,
les propriétés suivantes sont \omterm{def:b2:nvs:equivalent}{équivalentes}~:
\begin{enumerate}
  \item $u$ est \omterm{def:b2:metric:continuity}{continue}~;
  \item $u$ est \omterm{def:b2:metric:continuity}{continue} en $0$~;
  \item $u$ est bornée sur la boule unité fermée~: $\sup_{\norm x \leq
        1} \norm{u(x)} < \infty$~;
  \item il existe $C \geq 0$ tel que $\norm{u(x)} \leq C \norm x$ pour
        tout $x$~;
  \item $u$ est \omterm{def:b2:metric:continuity}{lipschitzienne}.
\end{enumerate}
Le plus petit tel $C$ est la \emph{norme d'opérateur}\index{norme d'opérateur}
$\vertiii{u} = \sup_{\norm x \leq 1}\norm{u(x)} = \sup_{x \neq 0}
\frac{\norm{u(x)}}{\norm x}$~; elle fait de l'espace $\mathcal{L}_c(E,
F)$ des applications linéaires \omterm{def:b2:metric:continuity}{continues} un espace normé, avec
\[
\vertiii{v \circ u} \leq \vertiii v\, \vertiii u .
\]
\end{theorem}

\begin{proof}
(1 $\Rightarrow$ 2) trivial. (2 $\Rightarrow$ 3)~: la continuité en $0$
avec $\varepsilon = 1$ fournit $\delta$ tel que $\norm x \leq \delta
\Rightarrow \norm{u(x)} \leq 1$~; l'homogénéité ramène alors tout $x$
vérifiant $\norm x \leq 1$ dans cette boule puis en revient~:
\[
\norm{u(x)} = \frac1\delta\,\norm{u(\delta x)} \leq
\frac1\delta ,
\]
car $\norm{\delta x} \leq \delta$. (3 $\Rightarrow$ 4)~:
pour $x \neq 0$, appliquer la borne à $\frac{x}{\norm x}$. (4
$\Rightarrow$ 5)~: $\norm{u(x) - u(y)} = \norm{u(x - y)} \leq
C\norm{x - y}$. (5 $\Rightarrow$ 1) connu.

Axiomes de \omterm{def:b2:nvs:norm}{norme} pour $\vertiii\cdot$~: l'homogénéité et la séparation
sont claires ($\vertiii u = 0$ force $u = 0$ sur la boule, donc
partout)~; l'inégalité triangulaire vient de $\norm{(u + v)(x)} \leq
\norm{u(x)} + \norm{v(x)}$. Sous-multiplicativité~: $\norm{v(u(x))}
\leq \vertiii v\,\norm{u(x)} \leq \vertiii v \vertiii u \norm x$.
\end{proof}

\begin{example}\label{ex:b2:nvs:operatornorms}
Sur $\bigl(C(\intcc{0}{1}), \norm\cdot_\infty\bigr)$~: l'évaluation
$f \mapsto f(0)$ a pour \omterm{thm:b2:nvs:continuouslinear}{norme d'opérateur} $1$~; l'intégration $f \mapsto
\int_0^1 f$ a pour \omterm{def:b2:nvs:norm}{norme} $1$~; l'application $f \mapsto \int_0^1 t f(t)\dd t$
a pour \omterm{def:b2:nvs:norm}{norme} $\int_0^1 t\,\dd t = \frac12$ (borne supérieure par
l'inégalité triangulaire pour les intégrales~; atteinte en $f \equiv 1$).
Mais la dérivation, de $(C^1, \norm\cdot_\infty)$ vers $(C^0,
\norm\cdot_\infty)$, n'est \emph{pas} \omterm{def:b2:metric:continuity}{continue}~: $\norm{\sin(nx)}_\infty
= 1$ tandis que la dérivée a pour \omterm{def:b2:nvs:norm}{norme} sup $n$. Linéaire n'implique pas
\omterm{def:b2:metric:continuity}{continue} en dimension infinie.
\end{example}

\begin{example}[Deux normes, deux verdicts sur une même suite]\label{ex:b2:nvs:sqrtnxn}
Sur $C(\intcc01)$, posons $g_n(x) = \sqrt{n}\,x^n$. Alors
\[
\norm{g_n}_1 = \frac{\sqrt n}{n + 1} \longrightarrow 0,
\qquad
\norm{g_n}_2^2 = \frac{n}{2n + 1} \longrightarrow \frac12,
\qquad
\norm{g_n}_\infty = \sqrt n \longrightarrow \infty :
\]
une seule suite, trois \omterm{def:b2:nvs:norm}{normes}, trois comportements --- convergence vers
zéro, absence de convergence (les \omterm{def:b2:nvs:norm}{normes} se stabilisent à $\frac1{\sqrt2}$
mais la limite simple est $0$), et explosion. La masse qui se concentre
près de $x = 1$ est invisible pour $\norm\cdot_1$, à demi visible pour
$\norm\cdot_2$, dominante pour $\norm\cdot_\infty$. En dimension infinie,
«~converge-t-elle~?~» n'est pas une question portant sur une suite~:
c'est une question portant sur une suite \emph{et} une \omterm{def:b2:nvs:norm}{norme}.
\end{example}

\begin{method}[Calculer une norme d'opérateur]\label{met:b2:nvs:opnorm}
Toujours en deux temps. \emph{Borne supérieure~:} estimer
$\norm{u(x)}$ par $C\norm x$ à l'aide d'inégalités triangulaires, de
Cauchy--Schwarz ou de bornes intégrales --- cela prouve $\vertiii u
\leq C$. \emph{Témoin~:} exhiber soit un $x_0 \neq 0$ précis vérifiant
$\norm{u(x_0)} = C\norm{x_0}$ (la borne est atteinte), soit une
suite de vecteurs unitaires $x_n$ avec $\norm{u(x_n)} \to C$ (la
borne est approchée). Les deux temps sont obligatoires~: une borne
supérieure seule ne donne que $\vertiii u \leq C$, un témoin seul que
$\vertiii u \geq C$. En dimension infinie, le témoin peut devoir être
une suite --- le supremum n'a pas à être atteint
(\cref{exo:b2:nvs:8}).
\end{method}

\begin{example}[Les opérateurs diagonaux voient toutes les normes de la même façon]\label{ex:b2:nvs:diagonalnorm}
Pour $D = \operatorname{diag}(d_1, \dots, d_n)$ sur $K^n$ muni de
\emph{l'une quelconque} des \omterm{def:b2:nvs:norm}{normes} $\norm\cdot_1, \norm\cdot_2,
\norm\cdot_\infty$~: de $\abs{d_ix_i} \leq
\bigl(\max_j\abs{d_j}\bigr)\abs{x_i}$ coordonnée par coordonnée,
$\norm{Dx} \leq \max_j\abs{d_j}\,\norm x$~; et $x = e_{j_0}$
(un indice maximisant) l'atteint. Donc $\vertiii D =
\max_j\abs{d_j}$ dans les trois cas~: pour les applications
diagonales, toutes les \omterm{def:b2:nvs:norm}{normes} raisonnables racontent la même
histoire, le plus grand facteur d'étirement. Toute la difficulté des
\omterm{thm:b2:nvs:continuouslinear}{normes d'opérateur} porte sur le comportement \emph{non} diagonal ---
c'est pourquoi les \omterm{def:b2:nvs:norm}{normes} adaptées du problème du week-end de ce
chapitre (question 22) fonctionnent en forçant d'abord une matrice à
devenir diagonale.
\end{example}

\begin{example}[Sommes de colonnes~: le jumeau en norme $1$ de l'\cref{exo:b2:nvs:4}]\label{ex:b2:nvs:columnsums}
Sur $(\R^n, \norm\cdot_1)$, la \omterm{thm:b2:nvs:continuouslinear}{norme d'opérateur} d'une matrice $A$ est
la plus grande somme de \emph{colonne} en valeur absolue. Déroulons la
méthode~: pour $\norm x_1 \leq 1$,
\[
\norm{Ax}_1 = \sum_i\Bigl|\sum_j a_{ij}x_j\Bigr|
\leq \sum_j \abs{x_j}\sum_i\abs{a_{ij}}
\leq \Bigl(\max_j\sum_i\abs{a_{ij}}\Bigr)\norm x_1 ,
\]
et la borne est atteinte en $x = e_{j_0}$ pour une colonne maximisante
$j_0$ --- le témoin le plus net qui soit. Ainsi pour $A =
\begin{psmallmatrix}1 & -2\\ 3 & 1\end{psmallmatrix}$~:
$\vertiii A_1 = \max(1 + 3,\ 2 + 1) = 4$, tandis que
$\vertiii A_\infty = 4$ également (lignes) --- une coïncidence ici, non
une loi~: transposez les coefficients de la matrice de façon
asymétrique et les deux \omterm{def:b2:nvs:norm}{normes} se séparent. Les lignes pour
$\norm\cdot_\infty$, les colonnes pour $\norm\cdot_1$~: le moyen
mnémotechnique est que les vecteurs unitaires de chaque \omterm{def:b2:nvs:norm}{norme} (motifs
de signes, resp.\ vecteurs de base) sélectionnent les sommes
correspondantes.
\end{example}

\begin{proposition}[Applications bilinéaires]\label{prop:b2:nvs:bilinear}
Une application bilinéaire $b \colon E \times F \to G$ est \omterm{def:b2:metric:continuity}{continue} si
et seulement si $\norm{b(x,y)} \leq C\norm x\,\norm y$ pour un certain
$C$~; elle est alors \omterm{def:b2:metric:continuity}{lipschitzienne} sur les parties bornées. (Même
schéma de preuve~; le produit $(u, v) \mapsto v \circ u$ et la
multiplication matricielle en sont les exemples clés.)
\end{proposition}

\begin{proof}
Si la borne est vérifiée~:
\[
b(x,y) - b(x_0,y_0) = b(x - x_0,\, y) + b(x_0,\, y - y_0),
\]
donc $\norm{b(x,y) - b(x_0,y_0)} \leq C\norm{x - x_0}\norm y +
C\norm{x_0}\norm{y - y_0}$~: continuité en $(x_0, y_0)$, et une
borne \omterm{def:b2:metric:continuity}{lipschitzienne} là où $\norm x, \norm y \leq R$. Réciproquement,
la continuité en $(0,0)$ fournit $\delta$ tel que $\norm{b(x,y)} \leq 1$
sur $\norm x, \norm y \leq \delta$~; on dilate les deux variables.
\end{proof}

\section{Dimension finie}

\begin{theorem}[Équivalence des normes en dimension finie]\label{thm:b2:nvs:finitedim}
Sur un espace de dimension finie, \emph{toutes les \omterm{def:b2:nvs:norm}{normes} sont
\omterm{def:b2:nvs:equivalent}{équivalentes}}.\index{équivalence des normes} Par conséquent, en
dimension finie~: la convergence, le caractère \omterm{def:b2:metric:topology}{ouvert}, la compacité, la
complétude sont des notions indépendantes de la \omterm{def:b2:nvs:norm}{norme}~; \omterm{def:b2:metric:compact}{compact} $=$
fermé et borné~; l'espace est \omterm{def:b2:metric:complete}{complet}~; et toute application linéaire
(ou multilinéaire) \emph{issue d'}un espace de dimension finie est
\omterm{def:b2:metric:continuity}{continue}.
\end{theorem}

\begin{proof}
Fixons une base et identifions $E \simeq K^n$~; il suffit de comparer
toute \omterm{def:b2:nvs:norm}{norme} $N$ à $\norm\cdot_\infty$.

\emph{Un sens relève de l'\omterm{def:b2:structures:algebra}{algèbre}~:} $N(x) = N(\sum x_i e_i) \leq
\sum \abs{x_i} N(e_i) \leq C \norm x_\infty$ avec $C = \sum N(e_i)$.
Ceci montre aussi que $N$ est \omterm{def:b2:metric:continuity}{continue} sur $(K^n, \norm\cdot_\infty)$
(elle est $C$-lipschitzienne~: $\abs{N(x) - N(y)} \leq N(x - y)$).

\emph{L'autre relève de la topologie~:} la sphère unité $S = \{x :
\norm{x}_\infty = 1\}$ est fermée et bornée dans $(K^n,
\norm\cdot_\infty)$, donc compacte
(\cref{thm:b2:metric:compactprops}~(2), valable pour $\C^n \simeq
\R^{2n}$). La fonction \omterm{def:b2:metric:continuity}{continue} $N$ atteint son minimum $c$ sur
$S$~; $c > 0$ car $N$ ne s'annule qu'en $0 \notin S$. L'homogénéité
propage la borne~: $N(x) \geq c \norm{x}_\infty$ pour tout $x$.

Conséquences~: tous les énoncés se ramènent à $(K^n,
\norm\cdot_\infty)$, où ils sont connus
(\cref{thm:b2:metric:rncomplete},
\cref{thm:b2:metric:compactprops})~; une application linéaire $u$
issue d'un $E$ de dimension finie vérifie $\norm{u(x)} \leq \sum\abs{x_i}
\norm{u(e_i)} \leq C'\norm{x}_\infty$~: la borne (4) du
\cref{thm:b2:nvs:continuouslinear}.
\end{proof}

\begin{corollary}\label{cor:b2:nvs:closedsubspace}
Tout sous-espace de dimension finie d'un espace normé quelconque est
fermé.
\end{corollary}

\begin{proof}
Il est \omterm{def:b2:metric:complete}{complet} pour la \omterm{def:b2:nvs:norm}{norme} induite
(\cref{thm:b2:nvs:finitedim}), et les parties complètes sont fermées
(\cref{def:b2:metric:complete}).
\end{proof}

\begin{example}[Une meilleure approximation calculée par symétrie]\label{ex:b2:nvs:bestaffine}
Dans $\bigl(C(\intcc{-1}{1}), \norm\cdot_\infty\bigr)$, à quelle
distance $f(x) = \abs x$ se trouve-t-elle du sous-espace (fermé, de
dimension deux) des fonctions affines $a + bx$~? Par symétrie,
remplacer $a + bx$ par $a - bx$ laisse $\norm{f - (a \pm bx)}_\infty$
inchangée, et le milieu $a$ fait au moins aussi bien (inégalité
triangulaire sur la moyenne)~: il suffit de considérer les constantes.
Pour une constante $a$~:
\[
\norm{\abs x - a}_\infty = \max\,(1 - a,\ a)
\geq \frac12 ,
\]
minimisée en $a = \frac12$~: la distance vaut $\frac12$, atteinte
par la constante $\frac12$. Remarquons la courbe d'erreur $\abs x -
\frac12$~: elle atteint $\pm\frac12$ alternativement en $x = -1, 0,
1$ --- trois extrema de signes alternés pour une meilleure
approximation issue d'une famille à deux paramètres. Ce schéma
d'\emph{équioscillation} n'est pas fortuit~; c'est la signature
d'optimalité que le problème du week-end de ce chapitre transforme en
théorème de Tchebychev.
\end{example}

\begin{example}[Sous-espaces fermés contre sous-espaces denses]\label{ex:b2:nvs:closeddense}
Dans $E = \bigl(C(\intcc{0}{1}), \norm\cdot_\infty\bigr)$~: chaque
$\R_n[X]$ (polynômes de degré $\leq n$, restreints à
$\intcc01$) est un sous-espace de dimension finie, donc \emph{fermé}
--- une limite uniforme de polynômes de degré $\leq n$ en est un. Mais
la réunion $\R[X]$ de tous ces sous-espaces est \emph{dense} dans $E$
(le théorème d'approximation de Weierstrass, démontré au
\cref{ch:b2:funcseq}), et des sous-espaces propres denses sont aussi
peu fermés qu'on puisse l'être. La morale~: la fermeture des
sous-espaces est un privilège de la dimension finie~; empiler des
étages fermés peut construire un gratte-ciel dense.
\end{example}

\begin{theorem}[Riesz]\label{thm:b2:nvs:riesz}
La boule unité fermée d'un espace normé $E$ est compacte \emph{si et
seulement si} $\dim E < \infty$.\index{théorème de Riesz}
\end{theorem}

\begin{proof}
Dimension finie~: fermée et bornée suffit
(\cref{thm:b2:nvs:finitedim}).

Réciproquement, supposons $\dim E = \infty$. \emph{Lemme de Riesz~:}
pour tout sous-espace fermé propre $F \subsetneq E$ et tout $\varepsilon \in
\intoo{0}{1}$, il existe un vecteur unitaire $x$ tel que $d(x, F) \geq 1 -
\varepsilon$. Preuve~: choisir $y \notin F$, poser $\delta = d(y, F) > 0$
($F$ fermé), choisir $f \in F$ tel que $\norm{y - f} \leq
\frac{\delta}{1 - \varepsilon}$, et poser $x = \frac{y - f}{\norm{y -
f}}$~: pour tout $g \in F$,
\[
\norm{x - g} = \frac{\norm{y - (f + \norm{y-f}\,g)}}{\norm{y - f}}
\geq \frac{\delta}{\norm{y-f}} \geq 1 - \varepsilon ,
\]
le numérateur étant une distance de $y$ à un point de $F$.

Construisons maintenant des vecteurs unitaires $x_1, x_2, \dots$ par
récurrence~: $F_k = \operatorname{Vect}(x_1, \dots, x_k)$ est de
dimension finie, donc fermé (\cref{cor:b2:nvs:closedsubspace}) et
propre~; le lemme de Riesz avec $\varepsilon = \frac12$ fournit un
vecteur unitaire $x_{k+1}$ tel que $d(x_{k+1}, F_k) \geq \frac12$. La
suite vérifie $\norm{x_p - x_q} \geq \frac12$ pour $p \neq q$~: aucune
sous-suite convergente --- la boule unité n'est pas compacte.
\end{proof}

\begin{example}[Riesz comme détecteur de dimension]\label{ex:b2:nvs:rieszdetector}
$C(\intcc01)$ est-il de dimension finie~? Riesz répond sans exhiber
aucune famille libre infinie explicite~: la suite $f_n(x) = x^n$
appartient à la boule unité fermée et vérifie, pour $m > n$,
$\norm{f_n - f_m}_\infty \geq f_n(x_0) - f_m(x_0) > 0$ en des points
bien choisis --- quantifié proprement dans le problème du week-end de
ce chapitre (question 16), où une sous-suite reste à distance mutuelle
$\geq \frac14$. Aucune sous-suite convergente, donc la boule n'est pas
compacte, donc $\dim C(\intcc01) = \infty$ d'après le
\cref{thm:b2:nvs:riesz}. La compacité de la boule unité est une
dichotomie parfaite~: elle a lieu en dimension finie, échoue en
dimension infinie, sans terrain intermédiaire --- la géométrie seule
lit le type de dimension.
\end{example}

\section{Espaces de Banach}

\begin{definition}\label{def:b2:nvs:banach}
Un \emph{espace de Banach}\index{espace de Banach} est un espace normé
\omterm{def:b2:metric:complete}{complet}. Exemples~: tout espace normé de dimension finie
(\cref{thm:b2:nvs:finitedim})~; $\bigl(C(\intcc{a}{b}),
\norm\cdot_\infty\bigr)$ (\cref{thm:b2:metric:rncomplete})~;
$\mathcal{L}_c(E, F)$ lorsque $F$ est de Banach (même schéma de preuve
que pour les fonctions \omterm{def:b2:metric:continuity}{continues}). Contre-exemple~: $\bigl(C(\intcc{0}{1}),
\norm\cdot_1\bigr)$ (\cref{exo:b2:nvs:7}).
\end{definition}

\begin{example}[La norme d'opérateur de l'intégration]\label{ex:b2:nvs:integrationnorm}
Sur $\bigl(C(\intcc01), \norm\cdot_\infty\bigr)$, posons $T(f)(x) =
\int_0^x f(t)\,\dd t$ (un endomorphisme~: $T(f)$ est \omterm{def:b2:metric:continuity}{continue}).
Déroulons la~\cref{met:b2:nvs:opnorm}. Borne supérieure~:
\[
\abs{T(f)(x)} \leq \int_0^x\abs f \leq x\,\norm f_\infty \leq
\norm f_\infty ,
\]
donc $\vertiii T \leq 1$. Témoin~: $f \equiv 1$ donne $T(f)(x) =
x$ et $\norm{T(f)}_\infty = 1 = \norm f_\infty$~: atteinte,
$\vertiii T = 1$. Mais notons $\vertiii{T^2} = \frac12 \neq
\vertiii T^2$~: en effet $T^2(f)(x) = \int_0^x(x - t)f(t)\dd t$ vérifie
$\abs{T^2(f)(x)} \leq \frac{x^2}2\norm f_\infty$, atteinte de nouveau
en $f \equiv 1$~; et plus généralement $\vertiii{T^n} = \frac1{n!}$ ---
la borne sous-multiplicative $\vertiii T^n = 1$ se trompe d'une
factorielle. C'est exactement le phénomène que l'astuce d'itération du
problème du week-end du~\cref{ch:b2:metric} convertit en résolubilité
globale des équations différentielles linéaires.
\end{example}

\begin{theorem}[Convergence absolue dans les espaces de Banach]\label{thm:b2:nvs:absoluteconvergence}
Dans un \omterm{def:b2:nvs:banach}{espace de Banach}, si $\sum \norm{u_n} < \infty$ alors $\sum u_n$
converge, et $\norm{\sum u_n} \leq \sum\norm{u_n}$. (La théorie
complète des séries dans les espaces normés fait l'objet du
\cref{ch:b2:series}.)
\end{theorem}

\begin{proof}
Sommes partielles $S_N$~: pour $q > p$, $\norm{S_q - S_p} \leq
\sum_{n=p+1}^{q}\norm{u_n}$, qui tend vers $0$ (critère de Cauchy
pour la série réelle des \omterm{def:b2:nvs:norm}{normes})~: $(S_N)$ est de Cauchy, donc
convergente. L'inégalité passe à la limite depuis l'inégalité
triangulaire finie.
\end{proof}

\begin{example}[Exponentielle de matrice, premier contact]\label{ex:b2:nvs:matrixexp}
L'\emph{exponentielle de matrice}\index{exponentielle de matrice}~:
$\mathcal{M}_n(K)$ munie de n'importe quelle \omterm{def:b2:nvs:norm}{norme} sous-multiplicative
($\vertiii{AB} \leq \vertiii A \vertiii B$) est de Banach (dimension
finie). Alors, pour tout $A$,
\[
\eu^A = \sum_{k=0}^{\infty} \frac{A^k}{k!}
\]
converge absolument ($\vertiii{A^k/k!} \leq \vertiii A^k /k!$,
sommable)~: bien définie. Le~\cref{ch:b2:diffeq} l'exploite
systématiquement.
\end{example}

\begin{example}[Une série de Neumann qui s'arrête]\label{ex:b2:nvs:neumannfinite}
Pour $A = \begin{psmallmatrix}0 & \frac12\\ 0 &
0\end{psmallmatrix}$~: $\vertiii A < 1$ dans toute \omterm{thm:b2:nvs:continuouslinear}{norme d'opérateur}
construite sur les \omterm{def:b2:nvs:norm}{normes} de l'\cref{ex:b2:nvs:examples}, et $A^2 = 0$,
donc la série géométrique s'effondre~:
\[
(I - A)^{-1} = \sum_{k \geq 0} A^k = I + A =
\begin{pmatrix}1 & \tfrac12\\ 0 & 1\end{pmatrix},
\]
vérifié par $(I - A)(I + A) = I - A^2 = I$. La nilpotence
tronque la série exactement comme elle a tronqué l'exponentielle au
\cref{ch:b2:reduction}~; et l'exemple calibre les attentes~: l'inverse
de Neumann est une série infinie en général, un polynôme précisément
lorsque la perturbation est nilpotente, et l'erreur après $N$ termes
est toujours majorée par le reste géométrique $\vertiii A^{N+1}/(1 -
\vertiii A)$.
\end{example}

\begin{example}[L'exponentielle d'un générateur de rotation]\label{ex:b2:nvs:rotationexp}
Prenons $A = \begin{psmallmatrix}0 & -\theta\\ \theta &
0\end{psmallmatrix}$. Alors $A^2 = -\theta^2 I$, donc les puissances
cyclent avec une période quatre, et la série se scinde en parties
paire et impaire~:
\[
\eu^{A} = \sum_{k}\frac{A^k}{k!}
= \Bigl(\sum_{j}\frac{(-1)^j\theta^{2j}}{(2j)!}\Bigr) I
+ \Bigl(\sum_{j}\frac{(-1)^j\theta^{2j+1}}{(2j+1)!}\Bigr)
\frac{A}{\theta}
= \begin{pmatrix}
\cos\theta & -\sin\theta\\
\sin\theta & \cos\theta
\end{pmatrix},
\]
tous les réarrangements étant autorisés par la convergence absolue. L'
exponentielle d'un générateur antisymétrique est une rotation ---
calculée ici purement à partir de la série, trois chapitres avant que
l'équation différentielle $x' = Ax$ (\cref{ch:b2:diffeq}) n'explique
\emph{pourquoi}~: $\eu^{tA}$ est un mouvement circulaire uniforme.
L'éclairage final~: les identités entre séries matricielles se
démontrent exactement comme les scalaires, dès qu'une \omterm{def:b2:nvs:norm}{norme}
sous-multiplicative certifie la convergence absolue.
\end{example}

\begin{example}[La norme sup signifie uniforme~: le dictionnaire]\label{ex:b2:nvs:uniformdictionary}
L'énoncé $\norm{f_n - f}_\infty \to 0$ \emph{est} la convergence
uniforme~: un seul nombre, $\sup_x\abs{f_n(x) - f(x)}$, majore
l'erreur en tout point simultanément. Le dictionnaire à l'œuvre sur
$f_n(x) = x^n$ sur $\intcc{0}{1}$~: simplement, $f_n
\to 0$ sur $\intco{0}{1}$ et $f_n(1) = 1$~; en \omterm{def:b2:nvs:norm}{norme},
$\norm{f_n - 0}_\infty = 1 \not\to 0$, et de fait la limite simple
est discontinue, donc hors de portée d'une limite en
$\norm\cdot_\infty$ dans $C(\intcc01)$ (qui est fermé
pour les limites uniformes, \cref{thm:b2:metric:rncomplete}). Sur
$\intcc{0}{a}$, $a < 1$~: $\norm{f_n}_\infty = a^n \to 0$ ---
la convergence uniforme rétablie en rétrécissant le domaine. Tout
énoncé de convergence du~\cref{ch:b2:funcseq} est un énoncé
sur cette unique \omterm{def:b2:nvs:norm}{norme}~; garder le dictionnaire en tête réduit de
moitié ce chapitre.
\end{example}

\begin{remark}[Pièges courants]\label{rem:b2:nvs:pitfalls}
(i) Une \omterm{thm:b2:nvs:continuouslinear}{norme d'opérateur} dépend des \emph{deux} \omterm{def:b2:nvs:norm}{normes} choisies~: la
même matrice a une $\vertiii\cdot_\infty$ donnée par les sommes de
lignes (\cref{exo:b2:nvs:4}) et une $\vertiii\cdot_1$ différente
(sommes de colonnes)~; parler de «~la~» \omterm{def:b2:nvs:norm}{norme} d'une matrice sans nommer
les \omterm{def:b2:nvs:norm}{normes} sous-jacentes n'a pas de sens. (ii) $\vertiii{AB} \leq
\vertiii A\,\vertiii B$ est une inégalité, en général stricte ---
les puissances peuvent décroître bien plus vite que ne le suggère la
borne $\vertiii A^k$, ce qui est tout l'intérêt des \omterm{def:b2:nvs:norm}{normes} adaptées (
problème du week-end de ce chapitre, question 22). (iii) «~Linéaire
implique \omterm{def:b2:metric:continuity}{continue}~» est un privilège de la dimension finie~: la
dérivation sur les polynômes est linéaire et non bornée
(\cref{ex:b2:nvs:operatornorms}). (iv) La convergence absolue de
$\sum u_n$ n'est utile que lorsque l'espace est \omterm{def:b2:metric:complete}{complet}
(l'\cref{exo:b2:series:9} construit le contre-exemple). (v) En
dimension infinie, un supremum sur la boule unité est un vrai
supremum~: ne pas supposer qu'il est atteint
(\cref{exo:b2:nvs:8}).
\end{remark}

\begin{remark}[Perspectives au sein de ce volume]\label{rem:b2:nvs:perspectives}
Trois rendez-vous sont désormais fixés. Avec le~\cref{ch:b2:series}~:
dans un \omterm{def:b2:nvs:banach}{espace de Banach}, les séries absolument convergentes
convergent, si bien que les séries géométrique et exponentielle
d'opérateurs deviennent des outils quotidiens --- inverser $I - A$,
définir $\eu^{A}$ (\cref{ex:b2:series:neumann}). Avec le
\cref{ch:b2:funcseq} et le~\cref{ch:b2:powerseries}~: la convergence
des suites de fonctions et des séries entières est la convergence dans
$\bigl(C, \norm\cdot_\infty\bigr)$ (\cref{ex:b2:nvs:uniformdictionary}),
et le rayon de convergence est un énoncé sur les séries géométriques
qui dominent. Avec le~\cref{ch:b2:fourier}~: les \omterm{def:b2:nvs:norm}{normes}
$\norm\cdot_2$ et $\norm\cdot_\infty$ sont réellement en désaccord sur
$C(\intcc{0}{1})$ (\cref{ex:b2:nvs:sqrtnxn}), ce qui explique
exactement pourquoi la convergence en moyenne quadratique des séries
de Fourier et la convergence uniforme sont deux théorèmes différents
avec deux prix différents.
\end{remark}

\begin{remark}[Où ce chapitre est utilisé]
Les \omterm{thm:b2:nvs:continuouslinear}{normes d'opérateur} et la série géométrique animent les arguments
de perturbation du~\cref{ch:b2:diffcalc} (théorème d'inversion locale)
et l'\omterm{ex:b2:nvs:matrixexp}{exponentielle de matrice} du~\cref{ch:b2:diffeq}~; l'\omterm{thm:b2:nvs:finitedim}{équivalence
des normes} autorise silencieusement tout argument «~choisissez votre
\omterm{def:b2:nvs:norm}{norme} préférée~» du~\cref{ch:b2:funcseq} et au-delà~; et la
distinction dimension finie/infinie du théorème de Riesz --- rendue
quantitative dans le problème du week-end de ce chapitre --- est la
raison pour laquelle le volume de troisième année a besoin de nouveaux
outils (convergence faible, Arzelà--Ascoli, projections dans les
espaces de Hilbert) là où ce volume pouvait encore extraire des
sous-suites convergentes.
\end{remark}

\section{Exercices}

\begin{exercise}[$\star$]\label{exo:b2:nvs:1}
Sur $\R^2$, dessiner les boules unités de $\norm\cdot_1$, $\norm\cdot_2$,
$\norm\cdot_\infty$, et démontrer les inégalités $\norm x_\infty
\leq \norm x_2 \leq \norm x_1 \leq 2\norm x_\infty$ avec les meilleures
constantes en dimension $2$.
\end{exercise}

\begin{exercise}[$\star$]\label{exo:b2:nvs:2}
$N(f) = \abs{f(0)} + \norm{f'}_\infty$ est-elle une \omterm{def:b2:nvs:norm}{norme} sur
$C^1(\intcc{0}{1})$~? La comparer à $\norm{f}_\infty$~: une
inégalité tient, l'autre échoue (l'exhiber).
\end{exercise}

\begin{exercise}[$\star$]\label{exo:b2:nvs:3}
Calculer la \omterm{thm:b2:nvs:continuouslinear}{norme d'opérateur} de $u(f) = \int_0^1 f(t)\,\eu^t\,\dd t$
sur $\bigl(C(\intcc{0}{1}), \norm\cdot_\infty\bigr) \to \R$, et du
décalage $S(x_1, x_2, \dots, x_n) = (x_2, \dots, x_n, 0)$ sur
$(K^n, \norm\cdot_\infty)$.
\end{exercise}

\begin{exercise}[$\star\star$]\label{exo:b2:nvs:4}
Sur $(\R^n, \norm\cdot_\infty)$, démontrer que la \omterm{thm:b2:nvs:continuouslinear}{norme d'opérateur}
d'une matrice $A$ vaut $\vertiii A_\infty = \max_i \sum_j \abs{a_{ij}}$
(la plus grande somme de ligne en valeur absolue). La calculer pour
$\begin{pmatrix} 1 & -2\\ 3 & 1\end{pmatrix}$.
\end{exercise}

\begin{exercise}[$\star\star$]\label{exo:b2:nvs:5}
Démontrer que $GL_n(K)$ est \omterm{def:b2:metric:topology}{ouvert} dans $\mathcal{M}_n(K)$ et que $A
\mapsto A^{-1}$ y est \omterm{def:b2:metric:continuity}{continue}. \emph{Indication~: pour l'ouverture, si
$\vertiii H < \frac{1}{\vertiii{A^{-1}}}$ alors $A + H = A(I +
A^{-1}H)$ avec $\vertiii{A^{-1}H} < 1$, et $I + B$ est inversible
pour $\vertiii B < 1$ par la série géométrique
(\cref{thm:b2:nvs:absoluteconvergence})~; pour la continuité, majorer
$(A+H)^{-1} - A^{-1}$ à l'aide de la même série.}
\end{exercise}

\begin{exercise}[$\star\star$]\label{exo:b2:nvs:6}
Soit $\varphi$ une forme linéaire sur un espace normé $E$. Démontrer que
$\varphi$ est \omterm{def:b2:metric:continuity}{continue} si et seulement si $\ker\varphi$ est fermé.
\emph{(Si $\ker\varphi$ est fermé et $\varphi \neq 0$, choisir $a$
tel que $\varphi(a) = 1$ et $r > 0$ tel que $B(a, r) \cap \ker\varphi =
\emptyset$~; en déduire $\abs{\varphi(h)} \leq \frac{1}{r}\norm h$ par un
argument d'homogénéité sur $a - \frac{h}{\varphi(h)}$.)}
\end{exercise}

\begin{exercise}[$\star\star$]\label{exo:b2:nvs:7}
Démontrer que $\bigl(C(\intcc{0}{1}), \norm\cdot_1\bigr)$ n'est pas
\omterm{def:b2:metric:complete}{complet}~: montrer que les fonctions $f_n$, rampes affines de $0$ à
$1$ sur $\bigl[\frac12 - \frac1n, \frac12\bigr]$ (valeur $0$
avant, $1$ après), forment une suite de Cauchy sans limite \omterm{def:b2:metric:continuity}{continue} en
$\norm\cdot_1$.
\end{exercise}

\begin{exercise}[$\star\star\star$]\label{exo:b2:nvs:8}
Sur $E = C(\intcc{0}{1})$ muni de $\norm\cdot_\infty$, considérer
\[
\varphi(f) = \sum_{n \geq 1} (-1)^n\, 2^{-n} f\bigl(\tfrac1n\bigr).
\]
Démontrer que $\varphi$ est une forme linéaire \omterm{def:b2:metric:continuity}{continue} bien définie avec
$\vertiii\varphi = 1$, mais que le supremum définissant
$\vertiii\varphi$ n'est \emph{pas atteint} sur la boule unité fermée.
\emph{(Borne supérieure~: inégalité triangulaire. \omterm{def:b2:nvs:norm}{Norme} $= 1$~: construire
une fonction \omterm{def:b2:metric:continuity}{continue} $f_K$ avec $\norm{f_K}_\infty \leq 1$ et $f_K(\frac1n)
= (-1)^n$ pour $n \leq K$ --- les points $\frac1n$ sont isolés les uns
des autres. Non-atteinte~: l'égalité forcerait $f(\frac1n) =
(-1)^n$ pour tout $n$, incompatible avec la continuité de $f$ en $0$
puisque $\frac1n \to 0$.)}
\end{exercise}

\begin{exercise}[$\star\star\star$]\label{exo:b2:nvs:9}
Soit $E$ un espace normé dans lequel la boule unité fermée est compacte.
Redémontrer, sans invoquer le~\cref{thm:b2:nvs:riesz}, que toute
suite bornée admet une sous-suite convergente, et démontrer que toute
forme linéaire sur $E$ est \omterm{def:b2:metric:continuity}{continue} si et seulement si $\dim E < \infty$.
\emph{(Pour la dimension infinie, construire une forme discontinue en
la définissant librement sur une suite normalisée libre et en la
prolongeant --- en admettant l'existence d'un supplémentaire
algébrique.)}
\end{exercise}

\begin{exercise}[$\star\star$]\label{exo:b2:nvs:10}
Sur $C(\intcc{0}{1})$, démontrer $\norm f_1 \leq \norm f_2 \leq \norm
f_\infty$ \emph{(Cauchy--Schwarz pour la première)}, et montrer à l'aide
de la famille $f_n(x) = x^n$ qu'aucune inégalité ne peut être inversée
à une constante près~: les trois \omterm{def:b2:nvs:norm}{normes} sont deux à deux non
\omterm{def:b2:nvs:equivalent}{équivalentes}.
\end{exercise}

\begin{exercise}[$\star\star$]\label{exo:b2:nvs:11}
(Distance à un hyperplan) Soit $\varphi$ une forme linéaire \omterm{def:b2:metric:continuity}{continue}
non nulle sur un espace normé $E$. Démontrer que
\[
d\bigl(x, \ker\varphi\bigr) =
\frac{\abs{\varphi(x)}}{\vertiii\varphi}
\qquad (x \in E),
\]
et vérifier sur l'\cref{exo:b2:nvs:8} que l'infimum peut n'être
atteint par aucun point de l'hyperplan.
\end{exercise}

\begin{exercise}[$\star\star\star$]\label{exo:b2:nvs:12}
Sur $E = \R[X]$ (tous les polynômes), posons $N_1(P) =
\sup_{\intcc{0}{1}}\abs P$ et $N_2(P) =
\sup_{\intcc{0}{2}}\abs P$. Montrer que $N_1 \leq N_2$ mais que
$N_1$ et $N_2$ ne sont \emph{pas} \omterm{def:b2:nvs:equivalent}{équivalentes}~; en déduire que
l'identité $(E, N_2) \to (E, N_1)$ est une bijection linéaire \omterm{def:b2:metric:continuity}{continue}
dont l'inverse est discontinu. Montrer enfin que $(E, N_1)$ n'est
pas \omterm{def:b2:metric:complete}{complet} \emph{(sommes partielles de Taylor de $\eu^x$)}. Ces trois
phénomènes sont impossibles en dimension finie --- dire pourquoi.
\end{exercise}

\section{Problème~: meilleure approximation et théorème de Tchebychev}

Avec quelle qualité une fonction peut-elle être approchée par des
polynômes d'un degré donné, et quel polynôme le fait le mieux~? Du
côté de l'existence, la réponse relève de ce chapitre~: la compacité
en dimension finie fait exister les meilleures approximations. Du côté
explicite, un cas non trivial peut se résoudre complètement à la main
--- parmi tous les polynômes \emph{unitaires} de degré $n$, celui de
plus petite \omterm{def:b2:nvs:norm}{norme} sup sur $\intcc{-1}{1}$ est le polynôme de Tchebychev
(normalisé), de \omterm{def:b2:nvs:norm}{norme} $2^{1-n}$~: le \emph{théorème extrémal de
Tchebychev}. Le problème démontre les deux côtés, puis mesure à quel
point la compacité échoue en dimension infinie~: la boule unité de
$C(\intcc{0}{1})$ contient des constellations infinies de points à
distance mutuelle $1$.

\begin{problem}[{Problème du week-end --- théorème extrémal de Tchebychev
et géométrie de la boule unité}]
\label{pb:b2:nvs:1}
Les \omterm{def:b2:nvs:norm}{normes} sans indice sont des \omterm{def:b2:nvs:norm}{normes} sup sur le segment indiqué.

\medskip
\noindent\textbf{Partie I --- Meilleure approximation dans les espaces
normés.}
\begin{enumerate}
  \item Soit $F$ un sous-espace de dimension finie d'un espace normé
        $E$ et $x \in E$. Démontrer que la distance $d(x,
        F) = \inf_{f \in F}\norm{x - f}$ est \emph{atteinte}
        \emph{(se ramener à une partie fermée bornée de $F$ et utiliser
        le~\cref{thm:b2:nvs:finitedim})}.
  \item Une \omterm{def:b2:nvs:norm}{norme} est \emph{strictement convexe} lorsque $\norm u = \norm v
        = 1$ et $u \neq v$ impliquent $\bigl\Vert\frac{u +
        v}2\bigr\Vert < 1$. Montrer que $\norm\cdot_2$ sur $\R^n$
        est strictement convexe \emph{(identité du parallélogramme)}, et
        que $\norm\cdot_1$ et $\norm\cdot_\infty$ ne le sont pas pour
        $n \geq 2$.
  \item Démontrer que pour une \omterm{def:b2:nvs:norm}{norme} strictement convexe, la meilleure
        approximation de la question 1 est \emph{unique}.
  \item Dans $(\R^2, \norm\cdot_\infty)$, calculer toutes les meilleures
        approximations de $x = (0, 1)$ par la droite $F =
        \operatorname{Vect}\bigl((1,0)\bigr)$~: un intervalle de
        minimiseurs.
  \item Dans $\bigl(C(\intcc{a}{b}), \norm\cdot_\infty\bigr)$,
        montrer que la meilleure approximation de $f$ par des
        \emph{constantes} est unique, égale à $c^* = \frac{\max f
        + \min f}{2}$, de distance $\frac{\max f - \min
        f}{2}$~; calculer les deux pour $f(x) = x^2$ sur
        $\intcc{0}{1}$.
\end{enumerate}

\medskip
\noindent\textbf{Partie II --- Polynômes de Tchebychev.}
\begin{enumerate}[resume]
  \item Montrer qu'il existe exactement un polynôme $T_n$ tel que
        $T_n(\cos\theta) = \cos n\theta$ pour tout $\theta$
        \emph{(récurrence $T_{n+1} = 2XT_n - T_{n-1}$ issue de la
        formule d'addition du cosinus)}, que $\deg T_n = n$, et que
        son coefficient dominant vaut $2^{n-1}$ pour $n \geq 1$.
  \item Montrer $\abs{T_n} \leq 1$ sur $\intcc{-1}{1}$, avec
        $T_n(\eta_k) = (-1)^k$ aux $n + 1$ points $\eta_k =
        \cos\frac{k\pi}{n}$ ($k = 0, \dots, n$), et que les
        racines de $T_n$ sont les $n$ points
        $\cos\frac{(2k-1)\pi}{2n}$, entrelacées avec les $\eta_k$.
  \item Calculer $T_2, T_3, T_4$, et vérifier l'alternance de
        $T_3$ en $\eta_0, \dots, \eta_3 = 1, \frac12, -\frac12,
        -1$ par évaluation directe.
  \item Pour $\abs x \geq 1$, démontrer
        \[
        T_n(x) = \frac{\bigl(x + \sqrt{x^2 - 1}\bigr)^n +
        \bigl(x - \sqrt{x^2 - 1}\bigr)^n}{2},
        \]
        et en déduire $T_n(x) \sim \frac12\bigl(x + \sqrt{x^2 -
        1}\bigr)^n \to \infty$ géométriquement pour $x > 1$ fixé.
  \item Démontrer la loi de composition $T_m \circ T_n = T_{mn}$
        \emph{(vérifier sur $\intcc{-1}{1}$ et invoquer la rigidité
        des polynômes)}.
\end{enumerate}

\medskip
\noindent\textbf{Partie III --- Théorème extrémal de Tchebychev.}
Posons $Q_n = 2^{1-n}T_n$ (unitaire, d'après la question 6).
\begin{enumerate}[resume]
  \item Soit $P$ unitaire de degré $n \geq 1$ vérifiant $\sup_{\intcc{-1}{1}}\abs P < 2^{1-n}$. En évaluant $D
        = Q_n - P$ aux points $\eta_k$ et en comptant les
        changements de signe, aboutir à une contradiction. Conclure~:
        \[
        \sup_{\intcc{-1}{1}}\abs P \;\geq\; 2^{1-n}
        \qquad\text{pour tout } P \text{ unitaire de degré } n.
        \]
  \item (Cas d'égalité) Supposons $\sup_{\intcc{-1}{1}}\abs P =
        2^{1-n}$ avec $P$ unitaire de degré $n$, et posons $D = Q_n
        - P \neq 0$. Montrer $(-1)^kD(\eta_k) \geq 0$ pour tout $k$~;
        montrer que chacun des $n$ intervalles
        $\intcc{\eta_{k}}{\eta_{k-1}}$ contient un zéro de $D$,
        et qu'un zéro partagé par deux intervalles consécutifs est
        un point \emph{intérieur} $\eta_k$ où $D' = 0$ également. En
        conclure que $D$ possède $n$ zéros comptés avec
        multiplicité, donc $D = 0$~: le minimiseur est exactement
        $Q_n$ --- le \emph{théorème extrémal de Tchebychev}.
  \item Reformuler le théorème comme une distance~: sur $\intcc{-1}{1}$,
        \[
        d_\infty\bigl(X^n,\ \R_{n-1}[X]\bigr) = 2^{1-n},
        \]
        avec pour unique meilleure approximation $X^n - Q_n$~; et montrer
        par la substitution affine $x = \frac{1+t}2$ que sur
        $\intcc{0}{1}$ la distance devient $2^{1-2n}$.
  \item (Nœuds d'interpolation optimaux) Pour $n$ nœuds $x_1, \dots,
        x_n \in \intcc{-1}{1}$, le polynôme nodal
        $\omega(x) = \prod_i(x - x_i)$ est unitaire de degré $n$.
        Déduire de la question 12 quel choix de nœuds minimise
        $\sup_{\intcc{-1}{1}}\abs\omega$, le facteur dépendant des
        nœuds dans la majoration classique de l'erreur
        d'interpolation, et donner la valeur minimale.
  \item Vérifier le cas $n = 2$ du théorème à la main (trouver
        directement $\inf_c \sup_{\intcc{-1}{1}}\abs{x^2 - c}$),
        et calculer numériquement la distance de la question 13 sur
        $\intcc{0}{1}$ pour $n = 10$. Que dit sa taille
        sur le graphe de $x^{10}$~?
\end{enumerate}

\medskip
\noindent\textbf{Partie IV --- La boule unité de
$C(\intcc{0}{1})$.}
\begin{enumerate}[resume]
  \item Soit $g_k(x) = x^{2^k}$. Montrer $\norm{g_k}_\infty = 1$ et
        $\norm{g_k - g_j}_\infty \geq \frac14$ pour $j > k$
        \emph{(évaluer au point où $x^{2^k} =
        \frac12$)}~: une suite bornée explicite sans sous-suite
        convergente --- la boule unité fermée n'est pas
        compacte, à la main.
  \item (Lemme de Riesz, affiné) Soit $F$ un
        sous-espace propre de \emph{dimension finie} d'un espace normé
        $E$. À l'aide de la question 1, produire un vecteur unitaire $x$
        tel que $d(x, F) = 1$ exactement --- et non seulement $\geq 1 -
        \varepsilon$ comme dans le lemme du~\cref{thm:b2:nvs:riesz}.
  \item En déduire~: dans tout espace normé de dimension infinie, il
        existe une suite de vecteurs unitaires de distances
        mutuelles $\geq 1$, et en redémontrer le théorème de Riesz.
  \item Dans $C(\intcc{0}{1})$, exhiber une telle constellation
        explicitement~: les fonctions triangles $h_n$ à support dans
        $\bigl[\frac1{n+1}, \frac1n\bigr]$ de valeur maximale $1$.
        Vérifier $\norm{h_n} = 1$, $\norm{h_n - h_m} = 1$ pour $n
        \neq m$, et remarquer que $h_n \to 0$ simplement mais pas
        uniformément.
  \item (La précompacité échoue) Montrer que la boule unité
        fermée de $C(\intcc{0}{1})$ ne peut être recouverte par un
        nombre fini de boules de rayon $\frac13$ \emph{(chaque telle
        boule contient au plus un $h_n$)} --- à comparer avec
        l'étape de précompacité dans la preuve du
        \cref{thm:b2:metric:borellebesgue}.
\end{enumerate}

\medskip
\noindent\textbf{Partie V --- Les \omterm{def:b2:nvs:norm}{normes} à l'œuvre sur les matrices, et
synthèse.}
\begin{enumerate}[resume]
  \item Démontrer que toute \omterm{def:b2:reduction:eigen}{valeur propre} $\lambda$ de $A \in
        \mathcal{M}_n(\C)$ vérifie $\abs\lambda \leq
        \vertiii A$ pour toute \omterm{thm:b2:nvs:continuouslinear}{norme d'opérateur}~; appliquer
        l'\cref{exo:b2:nvs:4} pour majorer les \omterm{def:b2:reduction:eigen}{valeurs propres} de
        $\begin{psmallmatrix}1 & -2\\ 3 & 1\end{psmallmatrix}$
        et comparer à leur module véritable.
  \item (\omterm{def:b2:nvs:norm}{Normes} adaptées) Soit $A$ \omterm{def:b2:reduction:diag}{diagonalisable}, $A =
        P\,\mathrm{diag}(\lambda_1, \dots, \lambda_n)\,P^{-1}$.
        Montrer que $N_P(x) = \norm{P^{-1}x}_\infty$ est une \omterm{def:b2:nvs:norm}{norme} dont la
        \omterm{thm:b2:nvs:continuouslinear}{norme d'opérateur} vérifie $\vertiii A_{N_P} =
        \max_i\abs{\lambda_i}$.
  \item En déduire~: pour $A$ \omterm{def:b2:reduction:diag}{diagonalisable}, $A^k \to 0$ si et
        seulement si toutes les \omterm{def:b2:reduction:eigen}{valeurs propres} vérifient $\abs{\lambda_i} < 1$ ---
        l'\omterm{thm:b2:nvs:finitedim}{équivalence des normes} rend la conclusion
        indépendante de la \omterm{def:b2:nvs:norm}{norme}. Vérifier sur $A =
        \frac14\begin{psmallmatrix}1 & 2\\ 2 &
        1\end{psmallmatrix}$.
  \item (Les constantes d'équivalence explosent) Sur $\R_n[X]$,
        comparer $N_c(P) = \max_k \abs{a_k}$ (coefficients) et
        $\norm{P}_{\intcc{0}{1}}$~: ce sont deux \omterm{def:b2:nvs:norm}{normes}, donc
        \omterm{def:b2:nvs:equivalent}{équivalentes} pour chaque $n$ fixé~; mais montrer, à l'aide du
        minimiseur unitaire de la question 13 sur $\intcc{0}{1}$, que la
        meilleure constante $C_n$ dans $N_c \leq
        C_n\norm\cdot_{\intcc{0}{1}}$ vérifie $C_n \geq
        2^{2n-1}$. Conclure en une phrase pourquoi «~toutes les \omterm{def:b2:nvs:norm}{normes}
        sont \omterm{def:b2:nvs:equivalent}{équivalentes}~» meurt en dimension infinie.
  \item (Synthèse) Une phrase pour chaque point~: où la compacité des
        boules de dimension finie a fonctionné (questions 1, 12)~; ce
        que gouverne la stricte convexité~; ce que la constellation
        des questions 18--19 détruit~; et comment la question 24
        quantifie l'échec. Nommer le sommet (le théorème extrémal de
        Tchebychev) et dire où la meilleure approximation
        trouve sa maison moderne (le théorème de projection sur les
        espaces de Hilbert, volume de troisième année, où la complétude
        remplace la compacité).
\end{enumerate}
\end{problem}
